% ===== SECTION 1: INTRODUCTION =====

\section{Introduction}

\paragraph{The rise of small-model agents.}
Large language models (LLMs) have demonstrated remarkable capabilities as autonomous agents, using tools to complete complex tasks \cite{yao2023react, schick2023toolformer}. However, their deployment comes with significant costs: GPT-4o costs \$2.50 per million input tokens, and frontier models require substantial GPU infrastructure. This has motivated a parallel trend toward \emph{small language models} (SLMs)---models with fewer than 10 billion parameters that can run on consumer hardware, edge devices, and laptops \cite{belcak2025small, erdogan2024tinyagent, zhang2024tinyllama}. Recent position papers argue that SLMs are not merely a cost-saving alternative but are inherently more suitable for many agentic workloads, where tasks are specialized and repetitive rather than open-ended \cite{belcak2025small}. These models offer compelling advantages: lower latency, reduced cost, privacy preservation through local execution, and energy efficiency.

In production systems, agents make many model calls per user request, and most of those calls are short, structured, and routine \cite{karmakar2026agentfloor}. Recent work has shown that small models can match frontier performance on structured tool-use tasks, with the strongest open-weight model achieving aggregate accuracy comparable to GPT-5 while being substantially cheaper and faster to run \cite{karmakar2026agentfloor}. This has led to growing deployment of SLM-based agents in real-world applications, from edge devices \cite{erdogan2024tinyagent} to resource-constrained business automation \cite{wang2026agentic, cho2026harness}.

\paragraph{The reliability blindspot.}
Rising accuracy scores on standard benchmarks suggest rapid progress. However, accuracy is a poor proxy for reliability in deployed systems. Reliability---whether an agent behaves consistently, withstands perturbations, recovers from failures, and operates safely---is a distinct capability axis that existing evaluations largely ignore \cite{rabanser2026science, gupta2026reliabilitybench, yagubyan2026consistency}.

Critically, the existing reliability literature has focused almost exclusively on large frontier models. ReliabilityBench \cite{gupta2026reliabilitybench} evaluates two large models (GPT-4o and Gemini 2.0 Flash). The Science of AI Agent Reliability \cite{rabanser2026science} evaluates 15 models, all of which are frontier models (GPT-5, Claude Opus, Gemini Pro). Claw-Eval \cite{claw2026eval} tests 14 frontier models. None of these studies include models under 10 billion parameters.

The small models that \emph{are} evaluated in the literature are tested for capability, not reliability. AgentFloor \cite{karmakar2026agentfloor} tests 16 open-weight models on a capability ladder, but does not measure consistency, robustness, or safety. The ``Right for Wrong Reasons'' study \cite{right2026wrong} evaluates reasoning quality in small models, but focuses on static question-answering rather than tool-use agent behavior. This leaves a critical gap: \textbf{we do not know how reliable small models are when deployed as autonomous agents}.

\paragraph{This work.}
We present the first comprehensive, multi-dimensional reliability evaluation of open-weight small language models as tool-using autonomous agents. Our contributions are fourfold:

\begin{enumerate}[leftmargin=*, nosep]
    \item \textbf{Reliability benchmark for small-model agents.} We introduce a benchmark suite of 31 tool-use tasks spanning 8 categories, designed to test the four dimensions of reliability: consistency, robustness, fault tolerance, and safety.
    
    \item \textbf{Comprehensive evaluation across 9 models.} We evaluate nine representative SLMs (1\,B to 9\,B parameters)---Llama 3.2 1B, Llama 3.2 3B, Phi-3.5-mini, DeepSeek-R1 7B, Qwen 2.5 7B, Qwen 2.5 Coder 7B, Mistral 7B, Llama 3.1 8B, and Gemma 2 9B---all running locally on a consumer GPU with 6\,GB VRAM.
    
    \item \textbf{Key empirical findings.} We show that: (a) neither capability nor reliability scales with parameter count---the correlation with 31-task accuracy is weak and non-significant ($r = 0.289$, $p = 0.451$), and with composite reliability it is weakly negative ($r = -0.179$); identical-size 7\,B models span 25.8--67.7\% accuracy and 37.8--85.0\% composite reliability, while a 1\,B model (Llama 3.2 1B) achieves higher composite reliability than both 8\,B and 9\,B models; (b) code-specialized training transfers selectively to reliability: Qwen 2.5 Coder 7B ties its general counterpart on capability (67.7\%) but dominates all reliability dimensions (85.0\% composite); (c) robustness under input perturbations shows the widest model-to-model variance (0--90\%, mean 47.8\%), safety failures affect all models at alarming rates (50--100\% failure to refuse harmful requests), and mean fault tolerance is just 33.3\%; (d) reasoning-distilled models (DeepSeek-R1 7B) rank last in capability (25.8\%), achieve 0\% accuracy on the reliability suite, and execute 2.6--9.9$\times$ slower than the rest of the cohort---chain-of-thought traces conflict with ReAct format requirements.
    
    \item \textbf{Open-source release.} We release all code, benchmark definitions, evaluation traces, and analysis scripts to support reproducible research.\footnote{https://github.com/Praansu/small-agent-reliability}
\end{enumerate}

Our findings have direct implications for practitioners deploying small-model agents: they reveal \emph{where} reliability breaks down, \emph{which} models to choose for specific deployment contexts, and \emph{how} to mitigate the most critical failure modes without scaling model size. As SLMs continue their rapid adoption in edge and privacy-sensitive applications, understanding their reliability characteristics is not merely an academic exercise---it is a practical necessity.
